\section{Introducción y objetivos}

La simulación de eventos discretos permite analizar sistemas de espera cuya
dinámica exacta es difícil de resolver analíticamente. Este proyecto aborda dos
situaciones con características distintas. La primera es un servidor único de
integración y pruebas continuas (CI/CD) que recibe trabajos de forma Poisson,
los procesa en orden FIFO y tiene cola prácticamente ilimitada; su naturaleza
M/M/1 permite validar el simulador contra fórmulas exactas. La segunda es un
Centro de Operaciones de Seguridad (SOC) con tres analistas en paralelo, dos
clases de alertas con prioridad no expropiativa, distribuciones de servicio no
exponenciales y un buffer físico estrictamente finito, donde el bloqueo de
alertas es una consecuencia operativa relevante.

Los objetivos son: (i) implementar un motor de eventos discretos reproducible y
separado de los modelos; (ii) modelar y validar el servidor M/M/1; (iii)
simular el SOC y estimar sus métricas con incertidumbre estadística; (iv)
responder los Escenarios A y B planteados por la gerencia; y (v) dejar un
proceso reproducible que genere resultados, figuras y tablas del informe desde
las mismas semillas.

El alcance se limita a los dos sistemas descritos y a los supuestos del
enunciado. No se modelan abandono de clientes, vacaciones de servidores,
reintentos ni cambios de capacidad durante la operación. Todas las unidades de
tiempo están en minutos.
